% Grado 1 -- generado por tools/build_tex.py a partir de raw/grado_01.txt
\section{Grado 1}\label{grado:1}

% PDF p. 8
\dba{1}{Identifica los usos de los números (como código, cardinal, medida, ordinal) y las operaciones (suma y resta) en contextos de juego, familiares, económicos, entre otros.}

\evidencias

\begin{itemize}
  \item Construye e interpreta representaciones pictóricas y diagramas para representar relaciones entre cantidades que se presentan en situaciones o fenómenos.
  \item Explica cómo y por qué es posible hacer una operación (suma o resta) en relación con los usos de los números y el contexto en el cual se presentan.
  \item Reconoce en sus actuaciones cotidianas posibilidades de uso de los números y las operaciones.
  \item Interpreta y resuelve problemas de juntar, quitar y completar, que involucren la cantidad de elementos de una colección o la medida de magnitudes como longitud, peso, capacidad y duración.
  \item Utiliza las operaciones (suma y resta) para representar el cambio en una cantidad.
\end{itemize}

\ejemplo

A partir de diversos materiales (recortes de periódico, revistas, facturas, noticias, etiquetas de productos alimenticios, la cuenta de servicios públicos, fotografías, placas de vehículos, números de documentos de identidad, entre otros) reconoce los números que aparecen allí. Identifica con cuáles de esos números:

\begin{itemize}
  \item Se puede conocer la cantidad de objetos de una colección.
  \item Pueden ordenar eventos u objetos.
  \item Pueden hacer operaciones. Propone preguntas que para ser resueltas requieren calcular una suma o una resta.
\end{itemize}
¿Se pueden sumar los números de una placa de un carro o moto? En caso afirmativo, ¿para qué sería útil ese dato? Es decir, ¿cuál es su interpretación?

% PDF p. 8
\dba{2}{Utiliza diferentes estrategias para contar, realizar operaciones (suma y resta ) y resolver problemas aditivos.}

\evidencias

\begin{itemize}
  \item Realiza conteos (de uno en uno, de dos en dos, etc.) iniciando en cualquier número.
  \item Determina la cantidad de elementos de una colección agrupándolos de 1 en 1, de 2 en 2, de 5 en 5.
  \item Describe y resuelve situaciones variadas con las operaciones de suma y resta en problemas cuya estructura puede ser a + b = ?, a + ? = c, o ? + b = c.
  \item Establece y argumenta conjeturas de los posibles resultados en una secuencia numérica.
  \item Utiliza las características del sistema decimal de numeración para crear estrategias de cálculo y estimación de sumas y restas
\end{itemize}

\ejemplo

Emplea una calculadora simple (o alguna aplicación que la simule) y explora el efecto que tiene el signo = (igual) a medida que se presiona varias veces después de digitar una suma o una resta.

\begin{itemize}
  \item Si se presiona 5 + 2 = = = ¿Cuál sería el resultado? q ¿Cuál sería el resultado si en la calculadora se presiona 4 + 3 = = = = = =?
  \item Describe las acciones que hace la calculadora. Si se digita el número 3 y luego se digita + 5 y se presiona la tecla igual diez veces, ¿cuáles números aparecerán en la calculadora cada vez que se digita un “igual”?
\end{itemize}

% PDF p. 9
\dba{3}{Utiliza las características posicionales del Sistema de Numeración Decimal (SND) para establecer relaciones entre cantidades y comparar números.}

\evidencias

\begin{itemize}
  \item Realiza composiciones y descomposiciones de números de dos dígitos en términos de la cantidad de “dieces” y de “unos” que los conforman.
  \item Encuentra parejas de números que al adicionarse dan como resultado otro número dado.
  \item Halla los números correspondientes a tener “diez más” o “diez menos” que una cantidad determinada.
  \item Emplea estrategias de cálculo como “el paso por el diez” para realizar adiciones o sustracciones.
\end{itemize}

\ejemplo

En una bolsa hay billetes de dos denominaciones \$1 y \$10. Con esos billetes realiza el siguiente juego con uno o varios compañeros de clase. “Ambos piensan en un número, y sacan de la bolsa los billetes que requieran para completar la cantidad representada por dicho número. Gana el juego quien forme la cantidad usando el menor número de billetes”.

utilizó el utilizó el utilizó el de \$10 de \$1 de \$10 1 y 3

\fig{g01_p09_2}{Tabla del juego de billetes: cantidad pensada por cada jugador, billetes que usó el primer, segundo y tercer jugador, y jugador ganador.}

y 5 billetes y 5 billetes de de de \$10 de \$1 de \$10 y 15 billetes y 5 billetes de de

Luego de realizar varias partidas, explica por qué el mismo número se puede obtener con diferente cantidad de billetes y explica por qué el jugador ganador logró ser exitoso.

Se dispone de máximo 9 billetes de \$10 y 30 billetes de \$1. Para formar la cantidad \$47, encuentra al menos 5 maneras distintas de formar la cantidad solicitada. Identifica la forma en la que se usan menos billetes y encuentra una regla para saberlo rápidamente.

% PDF p. 10
\dba{4}{Reconoce y compara atributos que pueden ser medidos en objetos y eventos (longitud, duración, rapidez, masa, peso, capacidad, cantidad de elementos de una colección, entre otros).}

\evidencias

\begin{itemize}
  \item Identifica atributos que se pueden medir en los objetos.
  \item Diferencia atributos medibles (longitud, masa, capacidad, duración, cantidad de elementos de una colección), en términos de los instrumentos y las unidades utilizadas para medirlos.
  \item Compara y ordena objetos de acuerdo con atributos como altura, peso, intensidades de color, entre otros y recorridos según la distancia de cada trayecto.
  \item Compara y ordena colecciones según la cantidad de elementos.
\end{itemize}

\ejemplo

A partir de una colección de objetos cotidianos de diferentes tamaños y pesos 1 , que sean comparables respecto a algún atributo, como una piña, un carro de juguete, una uva, un lápiz, una hoja de papel, una manzana, entre otros, los ordena respecto a su tamaño y su peso y discute sobre las condiciones de ubicación entre ellos. Establece diversos ordenamientos de acuerdo con alguna magnitud, por ejemplo, se toman cajas de diferentes tamaños y se llenan con materiales como plastilina, arroz y algodón de modo que en la caja más pequeña quede el mayor peso y argumenta las razones para dicho ordenamiento.

1 Término usado en el sentido informal, al tomar en cuenta que el concepto de masa se desarrolla en grados posteriores.

% PDF p. 10
\dba{5}{Realiza medición de longitudes, capacidades, peso, masa, entre otros, para ello utiliza instrumentos y unidades no estandarizadas y estandarizadas.}

\evidencias

\begin{itemize}
  \item Mide longitudes con diferentes instrumentos y expresa el resultado en unidades estandarizadas o no estandarizadas comunes.
  \item Compara objetos a partir de su longitud, masa, capacidad y duración de eventos.
  \item Toma decisiones a partir de las mediciones realizadas y de acuerdo con los requerimientos del problema.
\end{itemize}

\ejemplo

Se dispone de tiras o cuerdas de diferentes tamaños, como las que se presentan en la imagen.

\fig{g01_p10_1}{Tiras de colores de distintas longitudes (morada, roja, verde, azul, naranja, amarilla) para comparar y medir longitudes.}

Identifica: a) Las tiras de otros colores que pueden armar la b) El número de tiras que caben en c) La cantidad de tiras y

\fig{g01_p10_2}{Preguntas con tiras: cuántas tiras amarillas caben en la tira morada; comparación de cantidades de tiras.}

que se necesitan para medir el largo de un lápiz o un clip. ¿De cuál de las dos tiras se necesitan más ?, ¿Por qué?

d) Anticipa la cantidad de tiras amarillas que se necesitan para medir un objeto si conoce que para medirlo se requieren 3 tiras de color naranja.

% PDF p. 11
\dba{6}{Compara objetos del entorno y establece semejanzas y diferencias empleando características geométricas de las formas bidimensionales y tridimensionales (Curvo o recto, abierto o cerrado, plano o sólido, número de lados, número de caras, entre otros).}

\evidencias

\begin{itemize}
  \item Crea, compone y descompone formas bidimensionales y tridimensionales, para ello utiliza plastilina, papel, palitos, cajas, etc.
  \item Describe de forma verbal las cualidades y propiedades de un objeto relativas a su forma.
  \item Agrupa objetos de su entorno de acuerdo con las semejanzas y las diferencias en la forma y en el tamaño y explica el criterio que utiliza. Por ejemplo, si el objeto es redondo, si tiene puntas, entre otras características.
  \item Identifica objetos a partir de las descripciones verbales que hacen de sus características geométricas.
\end{itemize}

\ejemplo

A partir de la construcción de títeres con material reciclable y de la configuración de objetos como los que se muestran en las figuras siguientes, relaciona las formas y cuerpos geométricos y encuentra características similares y diferentes entre la forma de las figuras y los sólidos que los componen.

\fig{g01_p11_1}{Personaje (mariposa) construido con cuerpos geométricos (cilindro).}

\fig{g01_p11_2}{Personaje (pingüino) construido con cuerpos geométricos (esfera/cono).}

\fig{g01_p11_3}{Personaje (gato) construido con cuerpos geométricos (cilindros).}

\fig{g01_p11_4}{Personaje «Señora Bigotes» construido con un cono.}

\fig{g01_p11_5}{Personaje construido con un cubo.}

% PDF p. 11
\dba{7}{Describe y representa trayectorias y posiciones de objetos y personas para orientar a otros o a sí mismo en el espacio circundante.}

\evidencias

\begin{itemize}
  \item Utiliza representaciones como planos para ubicarse en el espacio.
  \item Toma decisiones a partir de la ubicación espacial.
  \item Dibuja recorridos, para ello considera los ángulos y la lateralidad.
  \item Compara distancias a partir de la observación del plano al estimar con pasos, baldosas, etc.
\end{itemize}

\ejemplo

En un plano que representa el salón de clases hay una marca (estrella roja) que indica el lugar donde se ocultó un objeto. Escribe instrucciones que se darían a alguien que está en la puerta del salón para que encuentre el objeto.

Determina si se pueden dar otras instrucciones para llegar al mismo sitio.

\fig{g01_p12_1}{Plano del salón de clase visto desde arriba: pupitres en filas, tablero, puerta y otros objetos, para describir ubicaciones.}

% PDF p. 12
\dba{8}{Describe cualitativamente situaciones para identificar el cambio y la variación usando gestos, dibujos, diagramas, medios gráficos y simbólicos.}

\evidencias

\begin{itemize}
  \item Identifica y nombra diferencias entre objetos o grupos de objetos.
  \item Comunica las características identificadas y justifica las diferencias que encuentra.
  \item Establece relaciones de dependencia entre magnitudes.
\end{itemize}

\ejemplo

Se tiene un dispensador para pasar agua de un recipiente a un vaso. Al servir agua en el vaso el volumen de los dos recipientes cambia, describe cuáles de las otras magnitudes cambian y explica la relación entre ambas. Elabora dibujos en diferentes momentos, cuando se llena 1,2,3 vasos etc.

% PDF p. 13
\dba{9}{Reconoce el signo igual como una equivalencia entre expresiones con sumas y restas.}

\evidencias

\begin{itemize}
  \item Propone números que satisfacen una igualdad con sumas y restas.
  \item Describe las características de los números que deben ubicarse en una ecuación de tal manera que satisfaga la igualdad.
  \item Argumenta sobre el uso de la propiedad transitiva en un conjunto de igualdades.
\end{itemize}

\ejemplo

Llena los espacios vacíos para que el resultado de la cadena azul y la cadena verde sean iguales. Indaga otras posibles soluciones.

\fig{g01_p13_1}{Niño frente a un tablero con operaciones incompletas: □ + 2 = □, 5 + □, 9 − 2 = □.}

Llena los espacios vacíos para que el resultado de la cadena roja sea mayor que el resultado de la cadena azul e indaga si hay otras soluciones.

% PDF p. 13
\dba{10}{Clasifica y organiza datos, los representa utilizando tablas de conteo y pictogramas sin escalas, y comunica los resultados obtenidos para responder preguntas sencillas.}

\evidencias

\begin{itemize}
  \item Identifica en fichas u objetos reales los valores de la variable en estudio.
  \item Organiza los datos en tablas de conteo y/o en pictogramas sin escala.
  \item Lee la información presentada en tablas de conteo y/o pictogramas sin escala (1 a 1).
  \item Comunica los resultados respondiendo preguntas tales como: ¿cuántos hay en total?, ¿cuántos hay de cada dato?, ¿cuál es el dato que más se repite?, ¿cuál es el dato que menos aparece?
\end{itemize}

\ejemplo

Como bienvenida al año escolar se les va a brindar a los alumnos de 1A un helado. Se les pide que informen sobre cuáles son los sabores de su preferencia. Los niños con la ayuda de la profesora hacen una consulta y presentan el siguiente gráfico con los resultados obtenidos:

\fig{g01_p13_2}{Pictograma «Preferencia de helados de los estudiantes de 1A»: número de paletas por sabor (fresa, uva, vainilla, chocolate, maracuyá).}

Escribe una frase sencilla para responder preguntas tales como: ¿cuántos niños prefieren el helado de sabor de fresa?, ¿cuántos de uva?, etc., y ¿cuántos helados se deben comprar en total?, ¿cuál es el sabor más escogido por los niños del curso 1A?

